In online learning, there is often a tension between simple natural greedy algorithms such as follow-the-leader ($\ftl$), which performs well for benign instances but has poor worst case guarantees, and conservative minimax optimal algorithms which provide optimal performance guarantees over the worst case instances. We introduce a new online learning algorithm -- titled \emph{Switching via Monotone Adapted Regret Traces} ($\bob$) -- that adapts to the data and achieves regret that is \emph{instance optimal}, i.e., simultaneously competitive on every input sequence compared to the performance of the follow-the-leader ($\ftl$) policy and the worst case guarantee of any other input policy $\wcalg$.  
The regret of the $\bob$ policy on \emph{any} input sequence is within a multiplicative factor $e/(e-1) \approx 1.58$ of the smaller of: 1) the regret obtained by $\ftl$ on the sequence, and 2) the upper bound on regret guaranteed by the given worst-case policy. This guarantee holds for every input sequence regardless of how it is generated, which implies the typical weaker `best-of-both-worlds' bounds. $\bob$ exemplifies the simple principle of {\em measured optimism}, assuming the best while preparing for the worst. It always begins by playing the opportunistic algorithm $\ftl$, and switches at most once during the time horizon to $\wcalg$ as needed if the estimated performance of $\ftl$ is poor. Our approach and results follow \bluet{by drawing on competitive analysis techniques from the ski-rental problem}. 
We complement our competitive ratio upper bounds with a fundamental lower bound showing that over all input sequences, no algorithm can get better than a $1.43$-factor of the minimum regret achieved by $\ftl$ and the minimax-optimal policy.
%We present a modification of $\bob$ that combines $\ftl$ with a ``small-loss" algorithm to achieve instance optimality between the regret of $\ftl$ and the small loss regret bound.



% In online learning, there is often a tension between simple natural greedy algorithms such as follow-the-leader ($\ftl$), which performs well for benign instances but has poor worst case guarantees, and conservative minimax optimal algorithms which provide optimal performance guarantees over the worst case instances. We introduce a new online learning algorithm -- titled \emph{Switching via Monotone Adapted Regret Traces} ($\bob$) -- that adapts to the data and achieves regret that is \emph{instance optimal}, i.e., simultaneously competitive on every input sequence compared to the performance of the follow-the-leader ($\ftl$) policy and the worst case guarantee of any other input policy $\wcalg$.  
% The regret of the $\bob$ policy on \emph{any} input sequence is within a multiplicative factor $e/(e-1) \approx 1.58$ of the smaller of: 1) the regret obtained by $\ftl$ on the sequence, and 2) the upper bound on regret guaranteed by the given worst-case policy. This guarantee holds for every input sequence regardless of how it is generated, which implies the typical weaker `best-of-both-worlds' bounds. $\bob$ exemplifies the simple principle of {\em measured optimism}, assuming the best while preparing for the worst. It always begins by playing the opportunistic algorithm $\ftl$, and switches at most once during the time horizon to $\wcalg$ as needed if the estimated performance of $\ftl$ is poor. Our approach and results follow from a reduction of instance optimal online learning to competitive analysis for the ski-rental problem. 
% We complement our competitive ratio upper bounds with a fundamental lower bound showing that over all input sequences, no algorithm can get better than a $1.43$-factor of the minimum regret achieved by $\ftl$ and the minimax-optimal policy.